\documentclass[11pt,a4paper]{article}
\usepackage[T1]{fontenc}
\usepackage[utf8]{inputenc}
\usepackage{amsmath,amssymb,amsthm}
\usepackage[margin=1in]{geometry}
\usepackage[colorlinks=true,linkcolor=blue,urlcolor=blue]{hyperref}
\theoremstyle{plain}
\newtheorem{theorem}{Theorem}
\newtheorem{lemma}[theorem]{Lemma}
\newtheorem{corollary}[theorem]{Corollary}
\theoremstyle{definition}
\newtheorem{remark}[theorem]{Remark}
\title{A proof of the conjectured recurrence for OEIS A308726}
\author{Adrian Perez Fontelles\\ \small Independent researcher}
\date{25 August 2026}
\begin{document}
\maketitle

\begin{abstract}
OEIS A308726 carries a conjectured linear recurrence with polynomial coefficients, of
order 5, contributed by R.~J.~Mathar. We prove it. The generating function of the
entry is algebraic, and the recurrence is equivalent to the vanishing of an explicit
differential operator applied to it: writing $\theta=x\,\frac{d}{dx}$, the sequence
$b(n)=\sum_{i}p_{i}(n)a(n-i)$ has generating function
$B(x)=\sum_{i}x^{i}\,(p_{i}(\theta+i)A)(x)$, so the conjecture holds for all $n>d$ exactly
when $B$ is a polynomial of degree $d$. Since $A$ is algebraic over $\mathbb{Q}(x)$, and every simple
algebraic extension of $\mathbb{Q}(x)$ is closed under $\theta$, the residual $B$ can be
computed in closed form. Here it equals $- 96 x^{4} + 24 x^{3} - 12 x^{2}$, a polynomial of degree $4$, which proves
the recurrence for every $n>4$.
\end{abstract}

\noindent\small 2020 Mathematics Subject Classification. 11B37, 05A15, 33F10.\normalsize

\section{The sequence and the conjecture}

OEIS A308726 is ``The number of permutations of length n and tier at most 1, that is, the number of permutations of length n sortable by two passes through a stack where outputting the longest prefix matching the identity permutation is prioritized''. It has offset $0$ and begins
\[
a(0),\dots,a(7)\;=\;1, 1, 2, 6, 24, 112, 556, 2811,\ \dots
\]
The entry gives the generating function
\begin{equation}\label{eq:gf}
A(x)\;=\;\sum_{n\ge 0} a(n)\,x^{n}\;=\;\frac{1}{\sqrt{1 - 4 x}} - \frac{\sqrt{2 \sqrt{1 - 4 x} - 1}}{2 x} + \frac{1}{x} - \frac{1}{2 x \sqrt{1 - 4 x}} ,
\end{equation}
and it carries the following comment:
\begin{quote}\small
\textbf{Conjecture:} D-finite with recurrence: 3*n*(n-1)*(n+1)*a(n) -n*(n-1)*(67*n-101)*a(n-1) +2*(n-1)*(286*n\^{}2-1112*n+1089)*a(n-2) +4*(-580*n\^{}3+4200*n\^{}2-10106*n+8049)*a(n-3) +24*(184*n\^{}3-1784*n\^{}2+5770*n-6221)*a(n-4) -96*(4*n-15)*(2*n-9)*(4*n-17)*a(n-5)=0. - \emph{R. J. Mathar}, Jan 27 2020
\end{quote}
As of the ``Last modified'' line on the live entry (2021-10-29, revision 33) the
statement is still recorded as a conjecture, and no proof appears among the entry's links,
formulas or programs.

Writing the conjectured relation in the form $\sum_{i=0}^{5}p_{i}(n)\,a(n-i)=0$,
the coefficient polynomials are
\[
p_{0}(n)=3 n \left(n - 1\right) \left(n + 1\right),\qquad p_{1}(n)=- n \left(n - 1\right) \left(67 n - 101\right),\qquad p_{2}(n)=2 \left(n - 1\right) \left(286 n^{2} - 1112 n + 1089\right),\qquad p_{3}(n)=- 4 \left(580 n^{3} - 4200 n^{2} + 10106 n - 8049\right),\qquad p_{4}(n)=24 \left(184 n^{3} - 1784 n^{2} + 5770 n - 6221\right),\qquad p_{5}(n)=- 96 \left(2 n - 9\right) \left(4 n - 17\right) \left(4 n - 15\right)
\]

\section{Recurrences as differential operators}

Let $A(x)=\sum_{n\ge0}a(n)x^{n}$ be the generating function of a sequence, and let
$\theta=x\,\dfrac{d}{dx}$, so that $\theta$ acts on $x^{m}$ as multiplication by $m$. For a
polynomial $p$ we write $p(\theta)$ for the corresponding differential operator, so that
$p(\theta)A=\sum_{m}p(m)a(m)x^{m}$.

\begin{lemma}\label{lem:transfer}
Let $p_{0},\dots,p_{r}$ be polynomials and put $b(n)=\sum_{i=0}^{r}p_{i}(n)\,a(n-i)$, with
the convention $a(m)=0$ for $m<0$. Then
\[
B(x):=\sum_{n\ge0}b(n)x^{n}=\sum_{i=0}^{r}x^{i}\,\bigl(p_{i}(\theta+i)A\bigr)(x).
\]
\end{lemma}

\begin{proof}
Fix $i$ and substitute $m=n-i$:
\[
\sum_{n\ge0}p_{i}(n)a(n-i)x^{n}=\sum_{m\ge0}p_{i}(m+i)a(m)x^{m+i}
=x^{i}\sum_{m\ge0}p_{i}(m+i)a(m)x^{m}=x^{i}\bigl(p_{i}(\theta+i)A\bigr)(x),
\]
the last step because $p_{i}(\theta+i)$ multiplies the coefficient of $x^{m}$ by
$p_{i}(m+i)$. Summing over $i$ gives the claim.
\end{proof}

\begin{corollary}\label{cor:criterion}
With the notation of Lemma~\ref{lem:transfer}, the recurrence
$\sum_{i}p_{i}(n)a(n-i)=0$ holds for every $n>d$ if and only if $B$ is a polynomial of
degree at most $d$.
\end{corollary}

\begin{proof}
$b(n)$ is the coefficient of $x^{n}$ in $B$, so $b(n)=0$ for all $n>d$ says exactly that
$B$ has no terms above $x^{d}$.
\end{proof}

This turns the conjecture into a closed-form identity. The point is that it can then be
decided exactly, with no truncation: $A$ is algebraic, so $B$ lies in the same field as
$A$, and one only has to recognise it.

\begin{lemma}\label{lem:field}
Let $P(x,y)\in\mathbb{Q}(x)[y]$ be squarefree of degree $m\ge1$ in $y$, let
$K=\mathbb{Q}(x)[y]/(P)$, and write $\alpha$ for the class of $y$ in $K$. Then $K$ is
closed under $\theta$. Explicitly,
\[
\alpha'=-\frac{P_{x}(x,\alpha)}{P_{y}(x,\alpha)},
\]
the quotient taken inside $K$, and for $u=\sum_{i<m}c_{i}(x)\,\alpha^{i}$ with
$c_{i}\in\mathbb{Q}(x)$,
\[
\theta(u)=x\Bigl(\sum_{i<m}c_{i}'(x)\,\alpha^{i}
 +\alpha'\sum_{i<m}i\,c_{i}(x)\,\alpha^{i-1}\Bigr),
\]
reduced modulo $P$.
\end{lemma}

\begin{proof}
$P(x,\alpha)=0$ holds identically in $x$, so differentiating gives
$P_{x}(x,\alpha)+P_{y}(x,\alpha)\,\alpha'=0$. Because $P$ is squarefree,
$\gcd(P,P_{y})=1$ in $\mathbb{Q}(x)[y]$, so $P_{y}(x,\alpha)$ is a unit of $K$ and the
displayed quotient is defined there; its inverse is produced by the extended Euclidean
algorithm. The formula for $\theta(u)$ is the chain rule, and reduction modulo $P$ returns
the value to the basis $1,\alpha,\dots,\alpha^{m-1}$. Every operation stays inside $K$.
\end{proof}

Consequently, if $A\in K$ then every $p_{i}(\theta+i)A$ is in $K$, and so is $B$. Writing
$B=\sum_{i<m}b_{i}(x)\,\alpha^{i}$ with $b_{i}\in\mathbb{Q}(x)$, the test of
Corollary~\ref{cor:criterion} becomes: $B$ is a polynomial precisely when $b_{i}=0$ for
every $i\ge1$ and $b_{0}$ has constant denominator. Both are decided by exact
cancellation of rational functions.

\section{The computation for A308726}

The generating function \eqref{eq:gf} is algebraic over $\mathbb{Q}(x)$;
its minimal polynomial is
\[
P(x,y) \;=\; 16 x^{5} y^{4} - 8 x^{4} y^{4} - 64 x^{4} y^{3} + 8 x^{4} y^{2} + x^{3} y^{4} + 32 x^{3} y^{3} + 94 x^{3} y^{2} - 48 x^{3} y + x^{3} - 4 x^{2} y^{3} - 48 x^{2} y^{2} - 28 x^{2} y + 52 x^{2} + 6 x y^{2} + 22 x y - 25 x - 3 y + 3 ,
\]
so \eqref{eq:gf} lies in $K=\mathbb{Q}(x)[y]/\bigl(P(x,y)\bigr)$ and is the root of that polynomial with the right
expansion at the origin. Applying Lemma~\ref{lem:transfer} to the coefficient polynomials
of Section~1 and reducing in $K$ by Lemma~\ref{lem:field}, every component $b_{i}$ of $B$
with $i\ge1$ cancels identically and the remaining rational part collapses to
\[
B(x)\;=\;- 96 x^{4} + 24 x^{3} - 12 x^{2} ,
\]
a polynomial of degree $4$.

\begin{theorem}
The conjectured recurrence
\[
3 n \left(n - 1\right) \left(n + 1\right)\,a(n) -  n \left(n - 1\right) \left(67 n - 101\right)\,a(n-1) + 2 \left(n - 1\right) \left(286 n^{2} - 1112 n + 1089\right)\,a(n-2) -  4 \left(580 n^{3} - 4200 n^{2} + 10106 n - 8049\right)\,a(n-3) + 24 \left(184 n^{3} - 1784 n^{2} + 5770 n - 6221\right)\,a(n-4) -  96 \left(2 n - 9\right) \left(4 n - 17\right) \left(4 n - 15\right)\,a(n-5) \;=\;0
\]
holds for every $n>4$.
\end{theorem}

\begin{proof}
Immediate from Corollary~\ref{cor:criterion} and the displayed value of $B$.
\end{proof}

\begin{remark}
The polynomial $B$ is not an error term: it records the boundary of the recurrence. For
$n\le4$ some of the shifted terms $a(n-i)$ fall outside the range of the sequence,
and $B$ collects exactly those contributions. Above that range the relation is exact.
\end{remark}

\section{Verification}

Every number above is an output of a script written separately from this note, and two
independent checks were run.

First, the generating function was checked against the entry rather than assumed: the
Taylor coefficients of \eqref{eq:gf} were expanded and compared term by term with the DATA
section of A308726, agreeing on all 13 terms compared.

Second, the recurrence itself was evaluated directly on the published terms, in exact
integer arithmetic and without reference to the generating function: for each $n$ in range
the sum $\sum_{i}p_{i}(n)a(n-i)$ was formed from the entry's own values and found to
vanish, for all 21 values of $n$ from $n=5$ upward. This is what
Corollary~\ref{cor:criterion} predicts from the degree of $B$, and it is a genuinely
different computation from the symbolic one: it never forms $B$, never differentiates, and
uses only integers.

The symbolic step is exact throughout. The residual is obtained by rational-function
cancellation in $K$, not by series truncation, so the identity $B=- 96 x^{4} + 24 x^{3} - 12 x^{2}$ is an
identity of functions and the theorem holds for all $n>4$ simultaneously.

\begin{thebibliography}{9}
\bibitem{oeis} The OEIS Foundation, \emph{The On-Line Encyclopedia of Integer Sequences},
\url{https://oeis.org}, 2026. Sequence A308726.
\bibitem{kp} M.~Kauers and P.~Paule, \emph{The Concrete Tetrahedron}, Springer, 2011,
Chapter 7 (holonomic closure properties and the translation between recurrences and
differential equations).
\bibitem{stanley} R.~P.~Stanley, \emph{Enumerative Combinatorics}, Volume 2, Cambridge
University Press, 1999, Chapter 6 (algebraic and D-finite generating functions).
\end{thebibliography}
\end{document}
